\lysh{14963}{4}

\begin{alist}
\item Αναζητούμε πραγματικούς αριθμούς $x$ των οποίων η απόσταση από το $4$ είναι
δύο μονάδες μεγαλύτερη από την απόστασή τους από το $2$. Δηλαδή
$d(x, 4) - d(x, 2) = 2$.

\item Έστω ότι τα σημεία $M, A, B$ αναπαριστούν στον άξονα των πραγματικών αριθμών,
τους αριθμούς $x, 2, 4$ αντίστοιχα, όπως φαίνεται στα παρακάτω σχήματα.
Παρατηρούμε ότι $d(2, 4) = (AB) = 2$.

Για κάθε αριθμό $x \in (-\infty, 2]$ είναι $d(x, 4) - d(x, 2) = (MB) - (MA) = (AB) = 2$.

(Το σχήμα παραλείπεται)

Για κάθε αριθμό $x \in (2, 4]$ είναι
\begin{gather*}
d(x, 4) - d(x, 2) < d(x, 4) + d(x, 2) = (MB) + (MA) = (AB) = 2\text{και άρα} \\
d(x, 4) - d(x, 2) \ne 2.
\end{gather*}

(Το σχήμα παραλείπεται)

Για κάθε αριθμό $x \in (4, +\infty)$ είναι $(MB) < (MA)$ οπότε $d(x, 4) < d(x, 2)$ δηλαδή
$d(x, 4) - d(x, 2) < 0$ και άρα $d(x, 4) - d(x, 2) \ne 2$.

(Το σχήμα παραλείπεται)

\item Όπως δείξαμε στο $\beta)$, αν για τον πραγματικό αριθμό $x$ ισχύει ότι
$|x - 4| - |x - 2| = 2$, τότε $x \in (-\infty, 2]$. (Η απόσταση $d(x, 2)$ είναι το $|x - 2|$ και αντίστοιχα το υπόλοιπο).

Το τριώνυμο $x^2 - 6x + 8$ έχει ρίζες τους αριθμούς $2$ και $4$ και γίνεται μη αρνητικό για
$x \in (-\infty, 2] \cup [4, +\infty)$.

Συνεπώς αν για τον πραγματικό αριθμό $x$ ισχύει ότι $|x - 4| - |x - 2| = 2$, τότε
$x \in (-\infty, 2]$ και $x^2 - 6x + 8 \ge 0$.
\end{alist}

